% The spoken body of the homily for the Twenty-ninth Sunday in Ordinary Time,
% Year A. Continuous preaching, ready to read aloud: no heading, no direction
% to the preacher, no citation inside the speech. The loci, the relation to
% the reviewed interpretations, the audience, the word count and the pace
% estimate stand in the terminal note, which no one reads out. The blank lines
% are the paragraph breaks a speaker recovers his place by, and the seven
% vertical spaces divide the speech into its eight movements: the trap and the
% strange word; the coin and the verb render; what bears God's image; the worn
% coin restamped; why true words are hypocrisy; Cyrus and Caesar's limit; the
% altar; and the one thing to do, with the return to the coin.
%
% Scripture is spoken in the Douay--Rheims, the leaf's study translation, and
% the speech names that Bible once, just before the first verse is quoted. No
% oration or antiphon is quoted in any Missal wording: the Collect, the Prayer
% over the Offerings and the Prayer after Communion are described, and the
% second Communion antiphon is named only as one of the two the Missal gives,
% through the verse of Mark it takes up. The agents of the scene are those the
% Gospel names.

In today's Gospel a coin is held up in the temple at Jerusalem, and Jesus asks
of it a question I would like us to carry home: Whose image is this?

The Pharisees send their disciples to him, with the Herodians, and they open
with praise. In the words of the old Douay Bible: ``Master, we know that thou
art a true speaker and teachest the way of God in truth. Neither carest thou
for any man: for thou dost not regard the person of men.'' Notice that every
word of that praise is true.

Then comes the question: ``Is it lawful to give tribute to Caesar, or not?''
Saint Augustine calls it a double noose: say yes, and Jesus maligns his own
nation; say no, and he is a rebel against Caesar.

And Jesus answers, ``Why do you tempt me, ye hypocrites?''

Hold on to that strange word. They have said nothing but the truth, and he
calls them hypocrites.

\bigskip

First, the coin. Jesus says neither yes nor no. He says, ``Shew me the coin of
the tribute.'' They bring him a denarius. He asks, ``Whose image and
inscription is this?'' And they say, ``Caesar's.''

They have answered their own question. Saint John Chrysostom says that Jesus
asked, not because he did not know, but so that they would be bound by their
own answers. The coin is already Caesar's.

Then Jesus changes one word. They had asked whether it was lawful to
\emph{give} tribute. He tells them to \emph{render} it. In the old
translation: ``Render therefore to Caesar the things that are Caesar's.'' To
render is to give back what is owed. Chrysostom heard the difference: ``this
is not to give but to render.''

Then he adds what nobody asked about: ``and to God, the things that are
God's.'' And he stops, without saying what those things are.

\bigskip

So what are they? Saint Jerome gives the plain answer: tithes, first-fruits,
offerings, the temple's dues beside the empire's. That is a fair reading of
the words. But Saint Augustine heard more in Jesus' question. If the coin
belongs to Caesar because it bears Caesar's image, what bears the image of
God? The answer stands on the first page of the Bible: ``Let us make man to
our image and likeness.'' God, Augustine says, is asking for his money, as
Caesar asks for his. And then he puts it in a single line: ``to Caesar his
money, to God yourselves.''

Here the Gospel changes our own question. We often ask what portion God is
owed: an hour on Sunday, something in the collection, a little of our time.
But on Augustine's reading the answer is larger than any portion. God is not
owed a second tax set beside Caesar's. Caesar may have the coin. God is owed
the one who carries it.

\bigskip

Augustine knew, though, that this image is not in mint condition. ``We are
God's money,'' he says: ``we have wandered away as coin from the treasury. The
impression that was stamped upon us has been rubbed out by our wandering.''
That is what sin does to the image of God in us: it wears the face smooth.

But Augustine goes on: ``He has come to refashion, for He it was that
fashioned us at first.'' Saint Ambrose, reading the same scene in Saint Luke,
says that Christ does not have Caesar's image, because he is the image of God.
The one who bears God's image without a flaw came to strike the worn coin
again. And he did it by giving himself back to the Father whole. In Saint
Mark's words, which one of the two Communion antiphons today takes up, the Son
of man came ``to give his life a redemption for many.''

For us that restamping began at the font. Our true good, Augustine says, is to
be stamped again by a new birth. And Ambrose reminds his people what they
promised there. We all renounce the world in words, he says, but not all of us
in affection, and we renounce it when we receive the sacraments. The contract
of faith, he says, is greater than the contract of money. And then he uses the
Gospel's own verb: render what you promised, while you are still in this body.

\bigskip

Now back to that strange word. Why are men who speak the truth called
hypocrites?

Chrysostom says they spoke the truth, ``not with an upright mind.'' Saint
Thomas Aquinas says that a hypocrite, properly speaking, has one thing in his
mouth and another in his heart. They praised him in order to use him.

Their words had given nothing away, and so, Chrysostom says, the rebuke could
not have been a guess. It was a sign that he knew their secret thoughts, a
proof of who he is. And he silenced them with gentleness.

This is where the Gospel turns toward us. The danger it shows is not the man
who refuses to pay. It is the one who pays correctly, says the right words,
and keeps his heart for himself: Ambrose's renunciation in words and not in
affection. A coin can be handed over with the hand alone. A self cannot. If
what God asks is ourselves, it can only be given with the heart.

At the beginning of this Mass the Collect asked that our will be always
devoted to God, and that we serve him with a sincere heart, a heart not
divided. The verse that greeted the Gospel comes from a sentence in which
Saint Paul wants the Philippians to be ``blameless and sincere children of
God.'' In the Latin Bible the word for sincere means single. Aquinas explains
it: God is simple, and we are single like his children when our intention is
toward one thing.

\bigskip

Caesar is not left with nothing: the questioners were carrying his coin, and
Jesus told them to render it. But the first reading shows where Caesar stands. The Lord
speaks to Cyrus, a pagan king of Persia, ``whose right hand I have taken hold
of,'' and twice he tells him, ``thou hast not known me.'' God's hand is on a
king who does not know him, so that all may know ``that there is none besides
me.'' Caesar's coin is minted in God's world.

So Caesar's claim is real, and it has a limit. If Caesar asks for anything
that harms godliness, Chrysostom says, ``such a thing is no longer Caesar's
tribute, but the devil's.'' Augustine tells of Christian soldiers under the
emperor Julian, who had abandoned the faith. When he ordered them to march,
they marched. When he ordered them to offer incense to idols, they refused.
``They distinguished their everlasting from their temporal master.'' They gave
the emperor their service, but not themselves.

\bigskip

In a few moments bread and wine will be brought forward, and the prayer over
them will call them God's own gifts before we have offered anything. What we
place on the altar was his already. The psalm told us how to come: ``Bring up
sacrifices, and come into his courts.'' Augustine tells us what present to
bring: ``Enter with an humble heart into the house of God, and thou hast
entered with a present.'' The bread and wine go up, and we go up with them.

After Communion we will ask that sharing in heavenly things carry us forward,
helped by the things of time and taught by the things of eternity. The coin
helps; it does not teach. What forms us is the one we receive: the image of
God, given for us and given to us, so that his image may be stamped in us
again.

\bigskip

So here is something to do this week. Each time a coin or a bill passes
through your hands, let it ask you Jesus' question: Whose image is this? Pay
what you owe, honestly; that is rendering to Caesar. Then ask what of yours
belongs to God, and give it with the heart behind the words.

Saint Paul remembers the Thessalonians' labour of charity, and Chrysostom says
of it: ``Merely to love is no labor at all. But to love genuinely is great
labor.'' Think of one person you have been treating correctly, but not with
your heart, and this week love that person genuinely. And when you come
forward today and answer Amen, let your words and your heart say it together.

A coin wears down as it passes from hand to hand. With us it is the other way:
the more we give ourselves back to God, the clearer his image becomes. Saint
John promises where that ends: ``we shall be like to him: because we shall see
him as he is.'' On that day no one will need to ask whose image we bear.
